\chapter{DC Circuits and the Resistor Color Code}
\label{ch:dc-circuits}

\section{Series Circuits}

Components are in \textbf{series} when they are connected end-to-end so
that the same current flows through each of them in turn, with no
alternate path.

\begin{center}
\begin{circuitikz}[scale=1]
\draw (0,0) to[battery1, l=$U$] (0,2)
      to[R, l=$R_1$] (2.5,2)
      to[R, l=$R_2$] (5,2)
      to[R, l=$R_3$] (5,0)
      -- (0,0);
\end{circuitikz}
\end{center}

\begin{keyidea}[Series circuit rules]
\begin{itemize}
\item Current is the \emph{same} through every component: $I_1 = I_2 = I_3 = I$.
\item Voltages \emph{add} to equal the source voltage:
$U = U_1 + U_2 + U_3$.
\item Resistances \emph{add}: $R_{\text{total}} = R_1 + R_2 + R_3 + \cdots$
\end{itemize}
\end{keyidea}

Because the same current flows through every series resistor, and
$U_n = IR_n$ for each one, the voltage across any single resistor in a
series chain is a fraction of the source voltage proportional to that
resistor's share of the total resistance --- the \textbf{voltage divider}
rule:
\[
U_n = U \times \frac{R_n}{R_{\text{total}}}.
\]

\begin{examplebox}[Voltage divider]
A 12\,V source drives a series circuit of $R_1 = 1\,\text{k}\Omega$ and
$R_2 = 3\,\text{k}\Omega$. Find the voltage across $R_2$.

\emph{Solution.} $R_{\text{total}} = 4\,\text{k}\Omega$, so
\[
U_2 = 12\,\text{V} \times \frac{3\,\text{k}\Omega}{4\,\text{k}\Omega} = 9\,\text{V}.
\]
\end{examplebox}

\section{Parallel Circuits}

Components are in \textbf{parallel} when they share the same pair of
nodes, so the same voltage appears across each of them.

\begin{center}
\begin{circuitikz}[scale=1]
\draw (0,0) to[battery1, l=$U$] (0,2.5) -- (1,2.5);
\draw (1,2.5) to[R, l=$R_1$] (1,0) -- (0,0);
\draw (1,2.5) -- (2.5,2.5) to[R, l=$R_2$] (2.5,0) -- (1,0);
\draw (2.5,2.5) -- (4,2.5) to[R, l=$R_3$] (4,0) -- (2.5,0);
\end{circuitikz}
\end{center}

\begin{keyidea}[Parallel circuit rules]
\begin{itemize}
\item Voltage is the \emph{same} across every branch: $U_1 = U_2 = U_3 = U$.
\item Currents \emph{add} to equal the total supplied current:
$I = I_1 + I_2 + I_3$.
\item Conductances add, so resistances combine by the reciprocal rule:
\[
\frac{1}{R_{\text{total}}} = \frac{1}{R_1} + \frac{1}{R_2} + \frac{1}{R_3} + \cdots
\]
\end{itemize}
\end{keyidea}

For exactly two resistors in parallel, this simplifies to the frequently
used \textbf{product-over-sum} shortcut:
\[
R_{\text{total}} = \frac{R_1 R_2}{R_1 + R_2}.
\]
A useful sanity check: the parallel combination of any set of resistors is
always \emph{smaller} than the smallest individual resistor, because
adding a parallel path always gives current an easier route.

\begin{examplebox}[Parallel resistance]
Find the total resistance of a 100\,$\Omega$ and a 150\,$\Omega$ resistor
in parallel.

\emph{Solution.}
\[
R_{\text{total}} = \frac{100 \times 150}{100 + 150} = \frac{15000}{250} = 60\,\Omega.
\]
Notice this is less than 100\,$\Omega$, the smaller of the two, as expected.
\end{examplebox}

Currents split between parallel branches in inverse proportion to
resistance (the \textbf{current divider} rule); for two resistors,
\[
I_1 = I_{\text{total}} \times \frac{R_2}{R_1+R_2}, \qquad
I_2 = I_{\text{total}} \times \frac{R_1}{R_1+R_2}.
\]

\section{Series--Parallel Combinations}

Real circuits usually combine both arrangements. The systematic method is
to redraw the circuit in stages: identify a sub-group of resistors that
are purely in series or purely in parallel, collapse that sub-group into
a single equivalent resistor, and repeat until only one resistor remains.
Work from the ``furthest'' part of the circuit from the source inward.

\begin{examplebox}[Series--parallel reduction]
$R_1 = 2\,\text{k}\Omega$ is in series with the parallel combination of
$R_2 = 4\,\text{k}\Omega$ and $R_3 = 4\,\text{k}\Omega$, all driven by a
10\,V source. Find the total current drawn from the source.

\emph{Solution.} First collapse the parallel pair:
$R_{23} = (4\times4)/(4+4) = 2\,\text{k}\Omega$. This is now in series
with $R_1$: $R_{\text{total}} = 2 + 2 = 4\,\text{k}\Omega$. Then
$I = U/R_{\text{total}} = 10\,\text{V}/4\,\text{k}\Omega = 2.5\,\text{mA}$.
\end{examplebox}

\section{Kirchhoff's Laws}

The series and parallel rules above are special cases of two more general
laws, published by Gustav Kirchhoff in 1845, which apply to \emph{any}
circuit topology, no matter how tangled.

\begin{keyidea}[Kirchhoff's Current Law (KCL)]
The sum of currents entering a node equals the sum of currents leaving it
--- charge cannot accumulate at a node in steady state:
\[
\sum I_{\text{in}} = \sum I_{\text{out}}.
\]
\end{keyidea}

\begin{keyidea}[Kirchhoff's Voltage Law (KVL)]
The sum of voltage rises and drops around any closed loop is zero ---
energy gained from sources in a loop equals energy dissipated by
components in that loop:
\[
\sum U_{\text{loop}} = 0.
\]
\end{keyidea}

KCL and KVL are the foundation of general circuit analysis techniques
(mesh and nodal analysis) taught in full electrical-engineering courses;
for amateur radio purposes, it is enough to recognize that the series and
parallel rules of this chapter are direct consequences of these two laws,
and to be able to apply KVL/KCL qualitatively --- for example, to
recognize that the voltage across a shorted component must be zero, or
that current from a source must ultimately return to that source.

\section{Superposition and Thévenin Equivalents (Brief Introduction)}

Two further techniques appear in more advanced treatments and occasionally
in advanced-class exams:

\begin{itemize}
\item \textbf{Superposition}: in a linear circuit with multiple sources,
the response (current or voltage) at any point equals the sum of the
responses caused by each source acting alone, with all other independent
voltage sources replaced by short circuits and all other independent
current sources replaced by open circuits.
\item \textbf{Thévenin's theorem}: any two-terminal network of linear
sources and resistors can be replaced, from the point of view of an
external load, by a single equivalent voltage source $U_{\text{Th}}$ in
series with a single equivalent resistance $R_{\text{Th}}$. This is
extremely useful for analyzing how a source (such as a transmitter's
final amplifier) will drive different loads, and reappears in
Chapter~\ref{ch:transmission-lines} when discussing impedance matching.
\end{itemize}

\section{The Resistor Color Code}
\label{sec:color-code}

Small fixed resistors are usually too small to print numeric values on,
so their value and tolerance are encoded as a sequence of colored bands.
The most common scheme uses four bands: two significant digits, a
multiplier, and a tolerance.

\begin{center}
\renewcommand{\arraystretch}{1.15}
\begin{tabular}{@{}lccc@{}}
\toprule
Color & Digit & Multiplier & Tolerance \\
\midrule
\cellcolor{black!85}\textcolor{white}{Black} & 0 & $\times 10^{0}$ & --- \\
\cellcolor{brown!70!black}\textcolor{white}{Brown} & 1 & $\times 10^{1}$ & $\pm 1\%$ \\
\cellcolor{red}Red & 2 & $\times 10^{2}$ & $\pm 2\%$ \\
\cellcolor{orange}Orange & 3 & $\times 10^{3}$ & --- \\
\cellcolor{yellow}Yellow & 4 & $\times 10^{4}$ & --- \\
\cellcolor{green!70!black}\textcolor{white}{Green} & 5 & $\times 10^{5}$ & $\pm 0.5\%$ \\
\cellcolor{blue}\textcolor{white}{Blue} & 6 & $\times 10^{6}$ & $\pm 0.25\%$ \\
\cellcolor{violet}\textcolor{white}{Violet} & 7 & $\times 10^{7}$ & $\pm 0.1\%$ \\
\cellcolor{gray!60}Grey & 8 & $\times 10^{8}$ & --- \\
White & 9 & $\times 10^{9}$ & --- \\
\cellcolor{yellow!70!orange}Gold & --- & $\times 10^{-1}$ & $\pm 5\%$ \\
\cellcolor{gray!30}Silver & --- & $\times 10^{-2}$ & $\pm 10\%$ \\
None & --- & --- & $\pm 20\%$ \\
\bottomrule
\end{tabular}
\end{center}

For a \textbf{4-band resistor}, read bands 1--2 as significant digits,
band 3 as the power-of-ten multiplier, and band 4 as tolerance:
\[
\text{Value} = (\text{digit}_1 \, \text{digit}_2) \times 10^{\text{multiplier}} \,\Omega.
\]
A \textbf{5-band resistor} (more common for precision, $\pm1\%$ or
better, resistors) uses three significant digits before the multiplier
band.

\begin{examplebox}[Reading the color code]
A resistor shows the bands Brown, Black, Red, Gold. Find its value and
tolerance.

\emph{Solution.} Brown = 1, Black = 0, so the significant digits are
``10.'' Red is the multiplier, $\times10^2 = \times100$. So
$R = 10 \times 100 = 1000\,\Omega = 1\,\text{k}\Omega$. Gold indicates
$\pm5\%$ tolerance, so the resistor's true value lies between 950\,$\Omega$
and 1050\,$\Omega$.
\end{examplebox}

\begin{examtip}
A popular mnemonic for the digit sequence 0--9 (Black, Brown, Red,
Orange, Yellow, Green, Blue, Violet, Grey, White) is: \emph{``B}ad
\emph{B}oys \emph{R}ob \emph{O}ur \emph{Y}oung \emph{G}irls \emph{B}ut
\emph{V}iolet \emph{G}ives \emph{W}illingly.'' Use whichever mnemonic
sticks; what matters on an exam is speed and accuracy, not elegance.
\end{examtip}

\section{Preferred Values (E-series)}

Resistors (and capacitors) are not manufactured at every possible value;
they follow standardized \textbf{preferred value} series (E6, E12, E24,
E96, \dots) where the number indicates how many values are spaced
per decade. The E12 series, for instance, uses twelve values per decade
spaced so that the maximum error between any achievable value and a
desired value, given the series tolerance, stays small:
\[
1.0,\ 1.2,\ 1.5,\ 1.8,\ 2.2,\ 2.7,\ 3.3,\ 3.9,\ 4.7,\ 5.6,\ 6.8,\ 8.2 \ (\times 10^{n}).
\]
This is why component catalogs list ``4.7\,k$\Omega$'' and
``5.6\,k$\Omega$'' resistors but not ``5.0\,k$\Omega$'' --- the nearest
E12 values are used instead, and circuit designs are built around
whatever the nearest standard value achieves.

\begin{quickreview}
\item Series: same current, voltages add, resistances add.
\item Parallel: same voltage, currents add, resistances combine by
reciprocal sum (or product-over-sum for two resistors).
\item Kirchhoff's Current Law and Voltage Law are the general principles
underlying the series/parallel shortcuts.
\item Resistor color bands encode two or three significant digits, a
multiplier, and a tolerance.
\item Component values follow standardized E-series preferred value
tables, not arbitrary round numbers.
\end{quickreview}

\begin{practiceproblems}
\item Three resistors, 100\,$\Omega$, 220\,$\Omega$, and 330\,$\Omega$,
are in series across a 9\,V battery. Find the total resistance, the
current, and the voltage across the 220\,$\Omega$ resistor.
\item Find the equivalent resistance of 470\,$\Omega$, 470\,$\Omega$, and
1\,k$\Omega$ in parallel.
\item A resistor's bands read Yellow, Violet, Orange, Gold. Find its
value and tolerance range.
\item Explain, using KCL, why the total current entering a series--parallel
network must equal the current leaving it, even though internal branch
currents differ.
\item Two 100\,$\Omega$ resistors and one 220\,$\Omega$ resistor are
connected so that the 220\,$\Omega$ is in series with the parallel
combination of the two 100\,$\Omega$ units, across a 15\,V source. Find
the total current.
\end{practiceproblems}
